\documentclass[11pt]{article}
\usepackage[margin=1in]{geometry}
\usepackage{amsmath,amssymb,amsthm}
\usepackage[colorlinks=true,urlcolor=blue]{hyperref}
\newtheorem{theorem}{Theorem}
\newtheorem{lemma}{Lemma}
\newcommand{\Fp}{\mathbb F_p}
\newcommand{\one}{\mathbf1}
\newcommand{\Prob}{\mathbb P}
\newcommand{\E}{\mathbb E}
\newcommand{\KL}{\operatorname{D}}
\title{Balanced inputs and arbitrary targets\protect\\in the undirected regular configuration model}
\author{Odd-block RSBM research: independent mixing extension}
\date{November 15, 2026}
\begin{document}
\maketitle

\paragraph{Status and scope.}
This note gives the balanced-input part of the strong point-mixing
extension needed for odd block sizes. The graph is the undirected
configuration model, not the symmetrized permutation model. It extends
the counting interface of Huang~\cite[Proposition 2.2]{Huang} from target
zero to arbitrary targets. The proof uses exponential tilting to justify
a \emph{relative} moderate-deviation local limit estimate. All primes and
degrees are fixed. It is not by itself the full sum over every input
vector in M\'esz\'aros's Theorem 1.6~\cite{Meszaros}.

\begin{theorem}\label{thm:main}
Fix $d\ge3$, an odd prime $p$, and $K>0$. Let $C_n$ be the adjacency
matrix of the uniform pairing of $dn$ labelled half-edges, in $n$
labelled buckets of size $d$. Loops contribute two on the diagonal.
Along all integers with $dn$ even, uniformly over vectors $q\in\Fp^n$
whose label counts satisfy
\[
 \max_a|n_a-n/p|\le K\sqrt{n\log n},
\]
and over all $r\in\Fp^n$, one has
\[
 \Prob(C_nq=r)=
 \begin{cases}
 (1+o(1))p^{-n+1},&\sum_i r_i=d\sum_iq_i,\\
 0,&\text{otherwise}.
 \end{cases}
\]
There is no requirement that $n$ be even, and no requirement that
$p$ be coprime to $d$.
\end{theorem}

All sums of field elements in the theorem are in $\Fp$; integer counts
and their ordinary real logarithms below are separate objects.

\section{Exact arbitrary-target count}
For a pairing, let $M_{ab}$ count half-edges in buckets labelled $a$
whose partner bucket is labelled $b$. The admissible type matrices obey
\[
 M=M^T\ge0,\qquad M_{aa}\text{ even},\qquad
 \sum_bM_{ab}=dn_a.
\]
Write $\mathcal M$ for this finite set, and
\[
 W(M)=\prod_{a<b}M_{ab}!\prod_a
       \frac{M_{aa}!}{2^{M_{aa}/2}(M_{aa}/2)!},
 \qquad T=(dn-1)!!.
\]
The output prescribes the additional congruences
\[
 \Phi_a(M):=\sum_b bM_{ab}=S_a,
 \qquad S_a=\sum_{i:q_i=a}r_i\pmod p.
\]
For $t\in\Fp$, let $X^{(t)}\in\mathbb Z^p$ be the count vector of a
uniform ordered $d$-tuple of field elements conditional on its sum being
$t$. For each $a$, let $Z_a$ be the sum of independent $X^{(r_i)}$ over
vertices with $q_i=a$. Counting allowed labels on the $d$ ordered
half-edges at each vertex, and then pairing like labelled endpoints,
gives the exact identity
\begin{equation}\label{eq:exact}
 \Prob(C_nq=r)=\frac{p^{(d-1)n}}{T}
       \sum_{M\in\mathcal M}W(M)\prod_a\Prob(Z_a=M_a).
\end{equation}
An incompatible congruence makes the corresponding probability zero.
This is the direct arbitrary-target version of Huang's Proposition 2.2.

For comparison, let $Y_a$ be the count vector of $dn_a$ independent
uniform field labels, and put
\[
 B(M)=W(M)\prod_a\frac{(dn_a)!}{\prod_bM_{ab}!}.
\]
Every pairing has exactly one induced assignment of partner labels, so
\begin{equation}\label{eq:baseline}
 \sum_{M\in\mathcal M}B(M)=T,\qquad
 W(M)\prod_a\Prob(Y_a=M_a)=p^{-dn}B(M).
\end{equation}
Consequently $B(M)/T$ is exactly the distribution of the pairing's type.

\section{The relative row local limit estimate}
\begin{lemma}\label{lem:row}
Fix $p,d$. For any sequence $t_1,\ldots,t_N\in\Fp$, let
$Z=\sum_{i=1}^NX_i^{(t_i)}$, independently, and let $Y$ count $dN$
independent uniform labels. Uniformly for integer $m$ satisfying
\[
 \sum_bm_b=dN,\qquad
 \max_b|m_b-dN/p|=O(\sqrt N\log N),
\]
one has
\[
 \Prob(Z=m)=
 \begin{cases}
 (1+o(1))p\Prob(Y=m),&\sum_bbm_b=\sum_it_i\pmod p,\\
 0,&\text{otherwise}.
 \end{cases}
\]
The error is uniform over the entire target sequence.
\end{lemma}
\begin{proof}
For $d\ge3$, every one or two specified positions of the conditioned
$d$-tuple are independent uniform field elements. Hence every
$X^{(t)}$, and also the count vector of $d$ unconditioned uniform
labels, have mean $\mu=(d/p)\one$ and common covariance
\[
 \Sigma=(d/p)I-(d/p^2)\one\one^T.
\]
After dropping coordinate zero, the covariance is positive definite.
All distributions have fixed finite support.

The difference lattice of $X^{(t)}$ consists exactly of integer vectors
with total zero and weighted sum zero modulo $p$. To check the absence
of further character restrictions, suppose the phase
$\sum_j\theta_{a_j}$ is constant modulo $2\pi$ over every tuple with
$\sum_ja_j=t$. Comparing $(a,b)$ with $(a+u,b-u)$, with the remaining
$d-2$ coordinates chosen to complete the prescribed sum, gives
\[
 \theta_{a+u}-\theta_a=\theta_b-\theta_{b-u}.
\]
Thus every increment depends only on $u$, and
$\theta_b=\theta_0+2\pi\ell b/p$ modulo $2\pi$ for some $\ell$.
Characters separate sublattices, proving the claim. After deleting the
zero coordinate this lattice has index $p$ in $\mathbb Z^{p-1}$.
The unconditioned count vector has difference lattice
$\mathbb Z^{p-1}$. These facts do not use $p\nmid d$.

Here are details ensuring a relative estimate throughout the stated
moderate window. Use the $p-1$ retained coordinates. For the independent
array, let $\Psi(z)=\sum_i\log\E e^{z^TX_i^{(t_i)}}$.
All its derivatives through any fixed order are bounded by constants
times $N$ in a fixed real neighborhood, uniformly over the sequence.
Its linear and quadratic terms are the common $N\mu$ and $N\Sigma$.
For $u=m-N\mu=O(\sqrt N\log N)$, the inverse function theorem gives
a real tilt $z$ with $\nabla\Psi(z)=m$ and
\[
 z=N^{-1}\Sigma^{-1}u+O(\|u\|^2/N^2),\qquad
 \|z\|=O(\log N/\sqrt N).
\]
Under this tilt, the covariance sum is $N\Sigma+O(N\|z\|)$, and
\[
 \Psi(z)-z^Tm=-\frac1{2N}u^T\Sigma^{-1}u
                 +O(\log^3N/\sqrt N).
\]
The same expansion holds for the unconditioned counts, with the same
quadratic term.

Fourier inversion at the \emph{tilted mean} now gives a relative local
limit estimate: near each maximal-modulus character, set
$\theta=N^{-1/2}v$. The centered log characteristic function equals
$-\tfrac12v^T\Sigma v+O(N^{-1/2}\|v\|^3+\|z\|\|v\|^2)$.
Restricting to $\|v\|\le\log N$ makes this error tend uniformly to
zero. The remainder within a fixed character neighborhood is bounded
by $e^{-c\|v\|^2}$, and the part outside these neighborhoods is
$O(\rho^N)$ for some $\rho<1$, uniformly over the sequence and small
tilts. The latter follows because all finitely many distributions have
the same maximal-character set, and all tilted support weights stay
bounded below. Thus the common Gaussian determinant gives the same
main term for the two tilted laws, multiplied by their lattice indices.
The $p$ phases of the conditioned array agree precisely on its stated
congruence class. Its index is $p$, whereas the ordinary index is one.
Removing the real tilt proves the ratio in the lemma.
\end{proof}

\section{Uniformly negligible type tails}
Put $\nu_a=n_a/n$, $m=M/(dn)$, and use natural logarithms. There is a
constant $A=A(p,d)$ such that, for every fixed $q,r,M$,
\begin{equation}\label{eq:tail}
 \Prob(C_nq=r,\text{type}=M)
 \le(n+1)^A\exp\left\{-nH(\nu)
       -\frac{d-2}{2}n\KL(m\|\nu\nu^T)\right\}.
\end{equation}
For clarity, this bound needs no local limit theorem. Write
$\pi_a(b)=m_{ab}/\nu_a$ and evaluate the nonnegative row generating
polynomial at $z_b=\pi_a(b)^{(d-1)/d}$. Its value for prescribed sum
$t$ is
\[
 F_t(\pi)=\sum_{b_1+\cdots+b_d=t}
               \prod_j\pi(b_j)^{(d-1)/d}\le1.
\]
Indeed AM--GM bounds each product by
$d^{-1}\sum_j\prod_{i\ne j}\pi(b_i)$, whose sum over the constrained
tuple is one. Coefficient evaluation therefore bounds all allowed
vertex assignments by
\[
 \prod_{a,b}\pi_a(b)^{-((d-1)/d)M_{ab}}.
\]
Stirling's inequalities for $W(M)/T$ give, up to $(n+1)^A$, the exponent
\[
 n\left[\frac d2\sum_{ab}m_{ab}\log m_{ab}
       -(d-1)\sum_{ab}m_{ab}\log\frac{m_{ab}}{\nu_a}\right]
 =-nH(\nu)-\frac{d-2}{2}n\KL(m\|\nu\nu^T).
\]
Zero coordinates follow by continuity or by omitting zero exponents.
This proves \eqref{eq:tail}.

Let $\mathcal R_n$ contain the types with
\[
 \max_{a,b}|M_{ab}-dn\nu_a\nu_b|\le\sqrt n\log n.
\]
Balancedness gives $H(\nu)=\log p-O(\log n/n)$. Pinsker's inequality
and \eqref{eq:tail}, summed over at most $(dn+1)^{p^2}$ types, show
\begin{equation}\label{eq:tailsmall}
 \sup_r\Prob(C_nq=r,\ M\notin\mathcal R_n)=o(p^{-n}).
\end{equation}
The same factorial calculation for $B(M)/T$ gives a polynomial factor
times $\exp\{-\tfrac12dn\KL(m\|\nu\nu^T)\}$. Thus the baseline type
lies in $\mathcal R_n$ with probability $1-o(1)$. Both statements remain
true if the window is increased or decreased by any fixed constant.

\section{Congruence cosets and the normalization factor}
For every type, symmetry and the fixed row sums imply
\[
 \sum_a\Phi_a(M)=d\sum_bbn_b\pmod p.
\]
These are the only restrictions on the central type's congruence vector.
For $j\ne0$, the bounded integral symmetric matrix
\[
 Z_j=2(e_j-e_0)(e_j-e_0)^T
\]
has row sums zero and even diagonal, and
\[
 \Phi(Z_j)=2j(e_j-e_0).
\]
Since $p$ is odd, these vectors span the hyperplane of total zero.
Every difference between two allowed congruence vectors is therefore
realized by a bounded integral combination of the $Z_j$, with
coefficients in $\{0,\ldots,p-1\}$.

For any one of these bounded shifts $Z$, all entries stay positive in
the central window for large $n$. The explicit factorial expression
\[
 B(M)=\frac{\prod_a(dn_a)!}
 {\prod_{a<b}M_{ab}!\prod_a2^{M_{aa}/2}(M_{aa}/2)!}
\]
gives, uniformly there,
\begin{equation}\label{eq:smooth}
 \frac{B(M+Z)}{B(M)}=1+o(1).
\end{equation}
In fact the linear term is
\[
 -\frac12\sum_{ab}Z_{ab}\log(dn/p^2)=0.
\]
The remaining logarithmic error is $O(\log n/\sqrt n)$, because the central entries differ from
$dn/p^2$ by $O(\sqrt n\log n)$.
Using the slightly enlarged and shrunken windows to handle boundaries,
\eqref{eq:smooth} and baseline tail concentration show that all
$p^{p-1}$ allowed congruence cosets have baseline probability
\[
 p^{-(p-1)}+o(1),
\]
uniformly over the balanced input.

For $M\in\mathcal R_n$, every row lies within the moderate window of
Lemma~\ref{lem:row}. Apply that lemma to \eqref{eq:exact} and then use
\eqref{eq:baseline}. For an admissible $r$, the central part is
\[
 (1+o(1))p^{-n+p}
 \Prob_{\mathrm{baseline}}(M\in\mathcal R_n,\ \Phi(M)=S)
 =(1+o(1))p^{-n+1}.
\]
The omitted part is $o(p^{-n})$ by \eqref{eq:tailsmall}. If the global
sum condition fails, $\one^TC_n=d\one^T$ makes the event impossible.
This completes the proof of Theorem~\ref{thm:main}.

\paragraph{What this resolves.}
There are at most $p^n$ balanced inputs, so their total contribution to
the sum of maximum-atom errors is $o(1)$. The remaining inputs require
a separate argument, especially the shrinking neighborhood of the
uniform profile and the neighborhood of constant vectors. No full
strong-mixing theorem or odd-block planted-eigenvalue result should be
inferred from this note alone.

\begin{thebibliography}{9}
\bibitem{Huang} J.~Huang, Invertibility of adjacency matrices for random
$d$-regular graphs, arXiv:1807.06465v2 (2018), Proposition 2.2,
equations (2.11)--(2.14), and Section 4.1.
\url{https://arxiv.org/pdf/1807.06465v2}.
\bibitem{Meszaros} A.~M\'esz\'aros, The distribution of sandpile groups
of random regular graphs, arXiv:1806.03736v4 (2020), Theorem 1.6 and
Section 7. \url{https://arxiv.org/pdf/1806.03736v4}.
\end{thebibliography}
\end{document}
